% Fictional manuscript written for documentation and tests; errors are intentional.
\documentclass{article}
\usepackage{amsmath,amssymb,amsthm}
\newtheorem{lemma}{Lemma}
\title{Averaging on coordinate spaces}
\author{A. Author}
\date{}
\begin{document}
\maketitle
\section{The averaging map}
Let $n\geq 2$ and work over $\mathbb{Q}$. For $x=(x_1,\ldots,x_n)$ set
\[
  P(x)=\frac{x_1+\cdots+x_n}{n}(1,\ldots,1).
\]
\begin{lemma}\label{lem:constant-vectors}
The image of $P$ is the span of $(1,\ldots,1)$.
\end{lemma}
\begin{proof}
Every image vector is constant, and $P(1,\ldots,1)=(1,\ldots,1)$.
\end{proof}
\section{Kernel calculation}
The kernel of $P$ has dimension $n$.
Indeed, a vector is in the kernel precisely when its coordinates sum to zero.
\section{Further properties}
The image calculation follows from the lemma above.
The coordinates is equal after averaging.
\end{document}
